% Options for packages loaded elsewhere
\PassOptionsToPackage{unicode}{hyperref}
\PassOptionsToPackage{hyphens}{url}
\PassOptionsToPackage{dvipsnames,svgnames,x11names}{xcolor}
%
\documentclass[
]{article}
\usepackage{amsmath,amssymb}
\usepackage{iftex}
\ifPDFTeX
  \usepackage[T1]{fontenc}
  \usepackage[utf8]{inputenc}
  \usepackage{textcomp} % provide euro and other symbols
\else % if luatex or xetex
  \usepackage{unicode-math} % this also loads fontspec
  \defaultfontfeatures{Scale=MatchLowercase}
  \defaultfontfeatures[\rmfamily]{Ligatures=TeX,Scale=1}
\fi
\usepackage{lmodern}
\ifPDFTeX\else
  % xetex/luatex font selection
\fi
% Use upquote if available, for straight quotes in verbatim environments
\IfFileExists{upquote.sty}{\usepackage{upquote}}{}
\IfFileExists{microtype.sty}{% use microtype if available
  \usepackage[]{microtype}
  \UseMicrotypeSet[protrusion]{basicmath} % disable protrusion for tt fonts
}{}
\makeatletter
\@ifundefined{KOMAClassName}{% if non-KOMA class
  \IfFileExists{parskip.sty}{%
    \usepackage{parskip}
  }{% else
    \setlength{\parindent}{0pt}
    \setlength{\parskip}{6pt plus 2pt minus 1pt}}
}{% if KOMA class
  \KOMAoptions{parskip=half}}
\makeatother
\usepackage{xcolor}
\usepackage[margin=1in]{geometry}
\usepackage{listings}
\usepackage{minted}
\usepackage{etoolbox}
\newcommand{\passthrough}[1]{#1}
\lstset{defaultdialect=[5.3]Lua}
\lstset{defaultdialect=[x86masm]Assembler}
\definecolor{CodeBlockText}{HTML}{F8F8F2}
\usemintedstyle{monokai}
\AtBeginEnvironment{minted}{\color{CodeBlockText}}
\BeforeBeginEnvironment{minted}{\vspace{-0.25\baselineskip}}
\AfterEndEnvironment{minted}{\vspace{-0.35\baselineskip}}
\renewcommand{\theFancyVerbLine}{\textcolor{Tan}{\scriptsize\arabic{FancyVerbLine}}}
\setlength{\parskip}{4pt plus 1pt minus 1pt}
\setminted{
  bgcolor=Sepia,
  frame=single,
  framesep=3pt,
  framerule=0.4pt,
  formatcom=\color{CodeBlockText},
  linenos,
  numbersep=6pt,
  breaklines,
  breakanywhere,
  fontsize=\small,
  baselinestretch=0.96,
  tabsize=4,
  autogobble
}
\setlength{\emergencystretch}{3em} % prevent overfull lines
\providecommand{\tightlist}{%
  \setlength{\itemsep}{0pt}\setlength{\parskip}{0pt}}
\setcounter{secnumdepth}{-\maxdimen} % remove section numbering
\ifLuaTeX
  \usepackage{selnolig}  % disable illegal ligatures
\fi
\IfFileExists{bookmark.sty}{\usepackage{bookmark}}{\usepackage{hyperref}}
\IfFileExists{xurl.sty}{\usepackage{xurl}}{} % add URL line breaks if available
\urlstyle{same}
\hypersetup{
  colorlinks=true,
  linkcolor={blue},
  filecolor={Maroon},
  citecolor={Blue},
  urlcolor={blue},
  pdfcreator={LaTeX via pandoc}}

\title{Scaling Laws Project API Quickstart}
\author{}
\date{Draft for student distribution}

\begin{document}
\maketitle

{
\hypersetup{linkcolor=}
\setcounter{tocdepth}{3}
\tableofcontents
}
\hypertarget{scaling-laws-project-api-quickstart}{%
\section{Scaling Laws Project API
Quickstart}\label{scaling-laws-project-api-quickstart}}

Status: Draft for student distribution

This quickstart explains the public request and response format for the
course API: check budget, submit runs, poll results, organize data, and
submit the final configuration. Course staff will provide the API URL
and your personal API key before the project begins.

The English and Chinese companion Jupyter notebooks included with the
project release contain the same workflow as executable cells with
sample response objects. Use the notebook version when you want to run
or copy code directly.

\hypertarget{setup}{%
\subsection{1. Setup}\label{setup}}

Install the Python \passthrough{\lstinline!requests!} dependency:

\begin{minted}{bash}
python -m pip install requests
\end{minted}

Set the API URL and API key:

\begin{minted}{python}
import json
import os
import sys
import time
import requests

API_BASE_URL = os.environ["SCALING_API_BASE_URL"].rstrip("/")
API_KEY = os.environ["SCALING_API_KEY"]
HEADERS = {"Authorization": f"Bearer {API_KEY}"}


def api_request(method, path, **kwargs):
    response = requests.request(
        method,
        f"{API_BASE_URL}{path}",
        headers=HEADERS,
        **kwargs,
    )
    try:
        body = response.json()
    except ValueError:
        body = {"raw_response": response.text}

    if response.status_code == 409:
        return body
    if response.status_code < 200 or response.status_code >= 300:
        print(
            json.dumps(body, ensure_ascii=False, indent=2, sort_keys=True),
            file=sys.stderr,
        )
        raise SystemExit(1)
    return body
\end{minted}

Your API key is personal. Do not share it or include it in your report.

\hypertarget{check-budget}{%
\subsection{2. Check Budget}\label{check-budget}}

\begin{minted}{python}
budget = api_request("GET", "/budget")
budget
\end{minted}

Example response:

\begin{minted}{text}
{
  "student_id": "student-1",
  "total_seconds": 43200,
  "reserved_seconds": 0,
  "charged_seconds": 0,
  "remaining_seconds": 43200
}
\end{minted}

Budget is reported in seconds. \passthrough{\lstinline!43200!} seconds
equals 12 official accelerator-hours.

\hypertarget{create-a-small-configuration}{%
\subsection{3. Create a Small
Configuration}\label{create-a-small-configuration}}

Start with a small configuration to verify the workflow, correctness,
and stability before spending substantial budget on experiments.

\begin{minted}{python}
config = {
    "model": {
        "num_hidden_layers": 2,
        "hidden_size": 128,
    },
    "training": {
        "train_tokens": 4096,
        "num_evals": 1,
        "learning_rate": 3e-4,
    },
}
\end{minted}

The API fills documented defaults such as sequence length, batch size,
evaluation count, optimizer settings, and model seed. It checks
submissions early and rejects incompatible model dimensions,
divisibility errors, non-finite numeric values, requests that exceed
your remaining budget, and similar issues.

\hypertarget{supported-training-configuration-format}{%
\subsection{4. Supported Training Configuration
Format}\label{supported-training-configuration-format}}

The same training configuration object is used in two places:

\begin{itemize}
\tightlist
\item
  \passthrough{\lstinline!POST /submit!}: pass it as
  \passthrough{\lstinline!config!}
\item
  \passthrough{\lstinline!POST /final\_submission!}: pass it as
  \passthrough{\lstinline!training\_config!}
\end{itemize}

The public configuration has exactly two top-level objects:

\begin{minted}{python}
config = {
    "model": {
        # model fields
    },
    "training": {
        # training and optimizer fields
    },
}
\end{minted}

Use the documented fields below. Fields marked required must be present.
Fields marked optional may be omitted; the API fills the documented
default.

Required fields:

\begin{itemize}
\tightlist
\item
  \passthrough{\lstinline!model.num\_hidden\_layers!}: positive integer,
  number of Transformer decoder blocks. Increasing this at fixed width
  usually increases non-embedding parameters and training FLOPs roughly
  linearly, so it is one of the main model-size knobs for scaling-law
  sweeps.
\item
  \passthrough{\lstinline!model.hidden\_size!}: positive integer, width
  of the residual stream, token embeddings, attention projections, and
  MLP input/output. This is a strong architecture-size knob; attention
  and MLP parameter counts grow roughly quadratically with it at fixed
  depth. It must be compatible with the attention head layout described
  below.
\item
  \passthrough{\lstinline!training.train\_tokens!}: positive integer,
  total number of training tokens to consume from the fixed course data
  order. The worker trains on the first
  \passthrough{\lstinline!training.train\_tokens!} tokens implied by
  that deterministic stream; this field does not reshuffle data or
  define epochs. Larger values spend more optimizer work and more
  accelerator budget.
\item
  \passthrough{\lstinline!training.learning\_rate!}: positive finite
  float, peak learning rate. With the \passthrough{\lstinline!cosine!}
  schedule, training warms from zero to this value and then decays
  toward
  \passthrough{\lstinline!training.learning\_rate * training.final\_lr\_fraction!};
  with the \passthrough{\lstinline!constant!} schedule, the worker uses
  this value throughout training. Tune this together with model size and
  batch size: too large often produces numerical failure, while too
  small undertrains.
\end{itemize}

Optional model fields:

\begin{itemize}
\tightlist
\item
  \passthrough{\lstinline!model.attention\_bias!}: boolean, default
  \passthrough{\lstinline!false!}. Adds learned bias terms to attention
  projections when enabled. Leave it fixed across comparable experiments
  so parameter-count changes are attributable to the intended scaling
  variable.
\item
  \passthrough{\lstinline!model.head\_dim!}: positive integer, default
  \passthrough{\lstinline!model.hidden\_size / model.num\_attention\_heads!}.
  This is the width of each attention head. It is usually derived from
  \passthrough{\lstinline!model.hidden\_size!} and
  \passthrough{\lstinline!model.num\_attention\_heads!}; supply it only
  when you want the API to check an explicit head layout.
\item
  \passthrough{\lstinline!model.intermediate\_size!}: positive integer,
  default \passthrough{\lstinline!4 * model.hidden\_size!}. This is the
  hidden width of the feed-forward/SwiGLU block. Larger values increase
  non-embedding parameters and FLOPs substantially; common sweeps keep
  it as a fixed multiple of
  \passthrough{\lstinline!model.hidden\_size!}.
\item
  \passthrough{\lstinline!model.num\_attention\_heads!}: positive
  integer, default \passthrough{\lstinline!1!}. This is the number of
  query attention heads. At fixed
  \passthrough{\lstinline!model.hidden\_size!}, more heads mean a
  smaller \passthrough{\lstinline!model.head\_dim!}; keep
  \passthrough{\lstinline!model.head\_dim!} valid for the GPU attention
  backend.
\item
  \passthrough{\lstinline!model.num\_key\_value\_heads!}: positive
  integer, default
  \passthrough{\lstinline!model.num\_attention\_heads!}. This would
  control grouped-query attention, but the current Stanford reference
  runtime does not support grouped-query attention, so set it equal to
  \passthrough{\lstinline!model.num\_attention\_heads!}.
\item
  \passthrough{\lstinline!model.rms\_norm\_eps!}: positive finite float,
  default \passthrough{\lstinline!1e-6!}. Epsilon added in RMSNorm for
  numerical stability. The default is the reference-trainer value;
  changing it is usually an ablation rather than a scaling-law knob.
\item
  \passthrough{\lstinline!model.rope\_theta!}: positive integer, default
  \passthrough{\lstinline!1000000!}. Rotary-position embedding frequency
  base. Leave the default unless you are deliberately studying
  position-encoding behavior.
\item
  \passthrough{\lstinline!model.tie\_word\_embeddings!}: boolean,
  default \passthrough{\lstinline!false!}. When true, the output
  classifier reuses the input embedding matrix, reducing parameters.
  Keep this choice fixed across runs that you want to compare directly.
\item
  \passthrough{\lstinline!model.dtype!}: one of
  \passthrough{\lstinline!float32!}, \passthrough{\lstinline!bfloat16!};
  default \passthrough{\lstinline!bfloat16!}.
  \passthrough{\lstinline!bfloat16!} is the normal GPU-efficient
  training dtype. \passthrough{\lstinline!float32!} uses more memory and
  runtime and is mainly useful for controlled numerical comparisons.
\item
  \passthrough{\lstinline!model.vocab\_size!}: positive integer, default
  \passthrough{\lstinline!50432!}. This must match the tokenized data
  requirements. The current course deployment uses
  \passthrough{\lstinline!EleutherAI/gpt-neox-20b!} token IDs and
  requires at least \passthrough{\lstinline!50432!}; increasing it adds
  embedding/output parameters that usually do not represent useful model
  capacity for this assignment.
\end{itemize}

Optional training fields:

\begin{itemize}
\tightlist
\item
  \passthrough{\lstinline!training.sequence\_length!}: positive integer,
  default \passthrough{\lstinline!1024!}. Number of tokens in each
  training sequence/context window. It affects memory use and the number
  of tokens per optimizer step:
  \passthrough{\lstinline!training.sequence\_length * training.train\_batch\_size!}.
\item
  \passthrough{\lstinline!training.train\_batch\_size!}: positive
  integer, default \passthrough{\lstinline!1!}. Number of sequences in
  each optimizer step. At fixed
  \passthrough{\lstinline!training.train\_tokens!}, larger batches
  produce fewer optimizer updates and usually require more memory.
\item
  \passthrough{\lstinline!training.validation\_batch\_size!}: positive
  integer, default \passthrough{\lstinline!1!}. Number of validation
  sequences evaluated at once. It does not change the training
  objective, but it must fit memory and satisfy validation-token
  divisibility.
\item
  \passthrough{\lstinline!training.num\_evals!}: positive integer,
  default \passthrough{\lstinline!10!}. Number of training
  chunks/evaluation points. The worker trains for
  \passthrough{\lstinline!training.train\_tokens / training.num\_evals!}
  tokens per chunk and records one validation loss after each chunk; the
  final validation loss is the last entry. For tiny smoke tests with
  only a few optimizer steps, set this to \passthrough{\lstinline!1!}.
\item
  \passthrough{\lstinline!training.model\_seed!}: non-negative integer,
  default \passthrough{\lstinline!0!}. Controls model initialization
  only. The course data order is fixed, so changing this seed does not
  reshuffle the training stream.
\item
  \passthrough{\lstinline!training.optimizer!}: one of
  \passthrough{\lstinline!adamw!}, \passthrough{\lstinline!sgd!};
  default \passthrough{\lstinline!adamw!}. AdamW is the default
  reference optimizer. SGD is available for controlled comparisons but
  usually needs separate learning-rate tuning.
\item
  \passthrough{\lstinline!training.lr\_schedule!}: one of
  \passthrough{\lstinline!cosine!}, \passthrough{\lstinline!constant!};
  default \passthrough{\lstinline!cosine!}.
  \passthrough{\lstinline!cosine!} uses
  \passthrough{\lstinline!training.warmup\_fraction!} and
  \passthrough{\lstinline!training.final\_lr\_fraction!};
  \passthrough{\lstinline!constant!} keeps the learning rate fixed at
  \passthrough{\lstinline!training.learning\_rate!}.
\item
  \passthrough{\lstinline!training.weight\_decay!}: non-negative finite
  float, default \passthrough{\lstinline!0.0!}. Decoupled AdamW weight
  decay strength. It is ignored by SGD in the reference runtime. Common
  AdamW baselines often use small values such as
  \passthrough{\lstinline!0.01!}.
\item
  \passthrough{\lstinline!training.adam\_beta1!}: finite float in
  \passthrough{\lstinline![0, 1)!}, default
  \passthrough{\lstinline!0.9!}. AdamW first moment coefficient. Ignored
  by SGD.
\item
  \passthrough{\lstinline!training.adam\_beta2!}: finite float in
  \passthrough{\lstinline![0, 1)!}, default
  \passthrough{\lstinline!0.95!}. AdamW second moment coefficient.
  Ignored by SGD.
\item
  \passthrough{\lstinline!training.adam\_epsilon!}: positive finite
  float, default \passthrough{\lstinline!1e-8!}. AdamW denominator
  epsilon for numerical stability. Ignored by SGD.
\item
  \passthrough{\lstinline!training.warmup\_fraction!}: finite float in
  \passthrough{\lstinline![0, 1)!}, default
  \passthrough{\lstinline!0.0!}. Fraction of optimizer steps used to
  increase the learning rate from zero to
  \passthrough{\lstinline!training.learning\_rate!} under the cosine
  schedule.
\item
  \passthrough{\lstinline!training.final\_lr\_fraction!}: finite float
  in \passthrough{\lstinline![0, 1]!}, default
  \passthrough{\lstinline!0.0!}. Final learning-rate multiplier for the
  cosine schedule. For example, \passthrough{\lstinline!0.1!} ends at
  ten percent of the peak learning rate.
\item
  \passthrough{\lstinline!training.gradient\_clip\_norm!}: positive
  finite float, default \passthrough{\lstinline!1.0!}. Global
  gradient-norm clipping threshold. Lower values can stabilize spikes
  but may slow learning if they clip most updates.
\end{itemize}

Important validation rules:

\begin{itemize}
\tightlist
\item
  If \passthrough{\lstinline!model.head\_dim!} is supplied,
  \passthrough{\lstinline!model.hidden\_size!} must equal
  \passthrough{\lstinline!model.num\_attention\_heads * model.head\_dim!}.
  If \passthrough{\lstinline!model.head\_dim!} is omitted,
  \passthrough{\lstinline!model.hidden\_size!} must be divisible by
  \passthrough{\lstinline!model.num\_attention\_heads!}, and the API
  derives \passthrough{\lstinline!model.head\_dim!}.
\item
  \passthrough{\lstinline!model.num\_attention\_heads!} must be
  divisible by \passthrough{\lstinline!model.num\_key\_value\_heads!},
  and the current Stanford reference runtime requires
  \passthrough{\lstinline!model.num\_key\_value\_heads == model.num\_attention\_heads!}
  because grouped-query attention is not supported.
\item
  \passthrough{\lstinline!model.head\_dim!} must be
  \passthrough{\lstinline!<= 128!} and a multiple of
  \passthrough{\lstinline!8!} for the current GPU attention backend. For
  example, use \passthrough{\lstinline!model.num\_attention\_heads: 16!}
  with \passthrough{\lstinline!model.hidden\_size: 2048!} so
  \passthrough{\lstinline!model.head\_dim!} is
  \passthrough{\lstinline!128!}.
\item
  \passthrough{\lstinline!model.intermediate\_size!} must be greater
  than or equal to \passthrough{\lstinline!model.hidden\_size!}.
\item
  \passthrough{\lstinline!model.vocab\_size!} must be at least
  \passthrough{\lstinline!50432!} for the current
  \passthrough{\lstinline!EleutherAI/gpt-neox-20b!} tokenized dataset.
  Smaller vocabularies can make token IDs exceed the embedding table and
  are rejected.
\item
  When course staff run workers with
  \passthrough{\lstinline!worker\_gpu\_count > 1!}, the Stanford FSDP
  parameter sharding also requires
  \passthrough{\lstinline!model.hidden\_size!} and
  \passthrough{\lstinline!model.intermediate\_size!} to be divisible by
  \passthrough{\lstinline!worker\_gpu\_count!}. If
  \passthrough{\lstinline!model.tie\_word\_embeddings!} is
  \passthrough{\lstinline!false!},
  \passthrough{\lstinline!model.vocab\_size!} must also be divisible by
  \passthrough{\lstinline!worker\_gpu\_count!}.
\item
  \passthrough{\lstinline!training.train\_tokens!} must be divisible by
  \passthrough{\lstinline!training.sequence\_length * training.train\_batch\_size!}.
\item
  \passthrough{\lstinline!training.train\_tokens!} must be at most
  \passthrough{\lstinline!500000000000!}.
\item
  Total optimizer steps equal
  \passthrough{\lstinline!training.train\_tokens!} divided by
  \passthrough{\lstinline!training.sequence\_length * training.train\_batch\_size!}.
\item
  Total optimizer steps must be divisible by
  \passthrough{\lstinline!training.num\_evals!}.
\item
  Validation batches must fit the course validation-token setting;
  incompatible \passthrough{\lstinline!training.sequence\_length!} or
  \passthrough{\lstinline!training.validation\_batch\_size!} values are
  rejected.
\item
  When course staff run workers with
  \passthrough{\lstinline!worker\_gpu\_count > 1!},
  \passthrough{\lstinline!training.train\_batch\_size!} and
  \passthrough{\lstinline!training.validation\_batch\_size!} must be
  divisible by \passthrough{\lstinline!worker\_gpu\_count!} because the
  Stanford trainer shards the batch axis across the FSDP mesh.
\item
  All numeric fields must be finite, and fields documented as positive
  must be strictly greater than zero. In particular,
  \passthrough{\lstinline!training.learning\_rate!} must be positive and
  finite.
\item
  Unsupported top-level, \passthrough{\lstinline!model!}, or
  \passthrough{\lstinline!training!} fields are rejected, including
  legacy model field names. This helps catch misspelled field names
  before they silently change an experiment.
\end{itemize}

If a final submission uses a configuration that fails these checks, the
API returns \passthrough{\lstinline!invalid\_final\_submission!} and
does not store that final submission.

\hypertarget{submit-a-run}{%
\subsection{5. Submit a Run}\label{submit-a-run}}

\begin{minted}{python}
submit_result = api_request(
    "POST",
    "/submit",
    json={
        "config": config,
        "requested_runtime_seconds": 300,
    },
)
submit_result
\end{minted}

Example success response:

\begin{minted}{text}
{
  "experiment_id": "exp-000001",
  "status": "queued",
  "budget_reserved_seconds": 300,
  "resource_warnings": [
    "very_low_optimizer_steps"
  ],
  "resolved_config": {
    "parameter_count_estimate": 4260096,
    "tokens_per_optimizer_step": 1024,
    "total_optimizer_steps": 4,
    "resource_warnings": [
      "very_low_optimizer_steps"
    ]
  }
}
\end{minted}

Queued and running experiments reserve the full requested runtime. If a
run finishes early, unused reserved budget is returned according to
course rules. \passthrough{\lstinline!resource\_warnings!} are advisory
and do not reject the run.

Example duplicate response:

\begin{minted}{text}
{
  "error": "duplicate_config",
  "message": "An identical configuration was already submitted.",
  "experiment_id": "exp-000001"
}
\end{minted}

Duplicate detection is scoped to your API key and does not expose other
students' experiment IDs or configs.

\hypertarget{poll-a-run}{%
\subsection{6. Poll a Run}\label{poll-a-run}}

\begin{minted}{python}
experiment_id = submit_result["experiment_id"]
terminal_statuses = {"completed", "failed", "cancelled", "system_failed"}

while True:
    experiment = api_request("GET", f"/experiment/{experiment_id}")
    status = experiment["status"]
    print("status:", status)

    if status in terminal_statuses:
        break

    time.sleep(30)
\end{minted}

Example completed response:

\begin{minted}{text}
{
  "experiment_id": "exp-000001",
  "status": "completed",
  "validation_losses": [4.2, 3.7],
  "final_validation_loss": 3.7,
  "failure_reason": "",
  "used_runtime_seconds": 180,
  "completed_at": "2026-07-16T15:00:00Z",
  "failed_at": "",
  "resource_warnings": []
}
\end{minted}

For completed runs, use
\passthrough{\lstinline!final\_validation\_loss!} as the run result.
\passthrough{\lstinline!used\_runtime\_seconds!} is the billable runtime
under the course runtime rules.
\passthrough{\lstinline!validation\_losses!} is an ordered history of
validation evaluations during training. It is not a confidence interval
and not \passthrough{\lstinline![lower, upper]!}; for a completed run,
\passthrough{\lstinline!final\_validation\_loss!} equals the last value
in \passthrough{\lstinline!validation\_losses!}.

Example failed response:

\begin{minted}{text}
{
  "experiment_id": "exp-000002",
  "status": "failed",
  "validation_losses": [4.6, 4.1],
  "final_validation_loss": null,
  "failure_reason": "timeout",
  "used_runtime_seconds": 600,
  "completed_at": "",
  "failed_at": "2026-07-16T15:00:00Z",
  "resource_warnings": []
}
\end{minted}

Partial validation losses from failed runs are diagnostics for
debugging. They are not final validation losses and should not be
treated as completed-run results.

\hypertarget{list-and-save-results}{%
\subsection{7. List and Save Results}\label{list-and-save-results}}

\begin{minted}{python}
experiments = api_request("GET", "/experiments")["experiments"]

rows = []
for exp in experiments:
    rows.append(
        {
            "experiment_id": exp["experiment_id"],
            "status": exp["status"],
            "final_loss": (
                exp["final_validation_loss"]
                if exp["status"] == "completed"
                else None
            ),
            "validation_losses": exp["validation_losses"],
            "used_runtime_seconds": exp["used_runtime_seconds"],
        }
    )
\end{minted}

Keep an experiment log with each run's purpose, the scaling-law
hypothesis it tested, and how the result changed your next experimental
decision. Make sure your submitted configuration can train reliably;
final training crashes caused by configuration robustness issues make
the final score invalid.

\hypertarget{submit-the-final-configuration}{%
\subsection{8. Submit the Final
Configuration}\label{submit-the-final-configuration}}

Your final submission includes a training config, a predicted final
loss, and a prediction interval. Invalid configs detected by the API
return \passthrough{\lstinline!invalid\_final\_submission!} and are not
stored.

\begin{minted}{python}
final_submission = api_request(
    "POST",
    "/final_submission",
    json={
        "training_config": config,
        "predicted_final_loss": 3.25,
        "predicted_final_loss_lower": 3.18,
        "predicted_final_loss_upper": 3.36,
    },
)
final_submission
\end{minted}

Example response:

\begin{minted}{text}
{
  "student_id": "student-1",
  "training_config": {
    "model": {"num_hidden_layers": 2, "hidden_size": 128},
    "training": {"train_tokens": 4096, "num_evals": 1, "learning_rate": 0.0003}
  },
  "predicted_final_loss": 3.25,
  "predicted_final_loss_lower": 3.18,
  "predicted_final_loss_upper": 3.36,
  "updated_at": "2026-07-16T15:00:00Z",
  "frozen_at": "",
  "frozen_by": "",
  "freeze_reason": ""
}
\end{minted}

You may submit multiple times before the deadline. The latest valid
final submission before the deadline is used for final evaluation.

Confirm the current final configuration:

\begin{minted}{python}
final_submission = api_request("GET", "/final_submission")
final_submission
\end{minted}

After the deadline, this endpoint returns the frozen record with
\passthrough{\lstinline!frozen\_at!},
\passthrough{\lstinline!frozen\_by!}, and
\passthrough{\lstinline!freeze\_reason!} populated.

\hypertarget{common-errors}{%
\subsection{9. Common Errors}\label{common-errors}}

Invalid API key:

\begin{minted}{text}
{
  "error": "invalid_api_key",
  "message": "Invalid or missing API key."
}
\end{minted}

Invalid configuration:

\begin{minted}{text}
{
  "error": "invalid_config",
  "message": "model.hidden_size must be divisible by model.num_attention_heads"
}
\end{minted}

Missing final submission:

\begin{minted}{text}
{
  "error": "final_submission_not_found",
  "message": "Final submission not found."
}
\end{minted}

Contact course staff if you have further questions.

\end{document}
